% ═══════════════════════════════════════════════════════════════════
\chapter{Sprint 0 -- Initialisation et étude technique}
% ═══════════════════════════════════════════════════════════════════

\section{Introduction}

Dans la continuité du chapitre précédent, le \textbf{Sprint 0} constitue la phase de
fondation du projet InvisiThreat. Il a pour objectif de formaliser les bases nécessaires
à un développement maîtrisé : identification des acteurs, spécification des besoins
fonctionnels et non fonctionnels, structuration du Product Backlog, planification des
sprints et étude technique initiale portant sur les environnements, les technologies
retenues et l'architecture du système.

% ───────────────────────────────────────────────────────────────────
\section{Identification des acteurs}

Avant d'entamer la présentation des diagrammes, il convient d'identifier les acteurs
en interaction avec le système. En modélisation UML, un acteur désigne un rôle joué
par une entité externe — utilisateur humain, dispositif matériel ou système tiers —
qui communique directement avec le système étudié. Il peut consulter ou modifier
l'état du système en échangeant des messages porteurs de données.

L'identification préalable des acteurs permet de structurer les responsabilités,
de délimiter les périmètres d'accès et de justifier les mécanismes d'autorisation
et de traçabilité qui seront mis en place. InvisiThreat implique cinq acteurs
distincts, dont les rôles sont synthétisés dans le tableau~\ref{tab:acteurs}.

\begin{table}[H]
\centering
\caption{Acteurs du système InvisiThreat et leurs responsabilités}
\label{tab:acteurs}
\renewcommand{\arraystretch}{1.3}
\begin{tabularx}{\textwidth}{p{3cm} X}
\hline
\rowcolor{gray!15}
\textbf{Acteur} & \textbf{Rôle principal} \\
\hline
\textbf{Admin}
  & Gouvernance globale de la plateforme, gestion des comptes utilisateurs,
    configuration des intégrations CI/CD et supervision des politiques de sécurité. \\
\textbf{Développeur}
  & Déclenchement des scans sur ses projets rattachés, consultation des résultats
    et application des correctifs identifiés. \\
\textbf{Security Manager}
  & Pilotage de la posture sécurité, priorisation des alertes selon le risque,
    production et diffusion de rapports de sécurité. \\
\textbf{Viewer}
  & Accès en lecture seule aux tableaux de bord et aux rapports, sans pouvoir
    intervenir sur les configurations ou les états des vulnérabilités. \\
\textbf{Systèmes externes}
  & Pipelines CI/CD, dépôts Git et moteur OWASP~ZAP interagissant avec la
    plateforme via l'API REST pour déclencher des scans ou récupérer des résultats. \\
\hline
\end{tabularx}
\end{table}

% ───────────────────────────────────────────────────────────────────
\section{Analyse des besoins}

L'analyse des besoins d'InvisiThreat se situe à l'intersection de contraintes métier,
techniques et sécuritaires. Sur le plan métier, l'enjeu central est de garantir
une continuité de livraison logicielle sans compromettre la sécurité, tout en offrant
aux équipes une visibilité claire et exploitable sur le niveau de risque par projet.
Sur le plan technique, la solution doit s'intégrer à l'écosystème DevOps existant
(pipelines CI/CD, dépôts Git, environnements conteneurisés) et traiter des volumes
d'alertes croissants sans dégrader les temps de réponse.

Les besoins des équipes sécurité et développement convergent vers une même exigence :
disposer d'un processus structuré de gestion des vulnérabilités, fondé sur la
traçabilité, la priorisation et l'automatisation. Cela suppose un modèle de données
commun pour les alertes, une historisation par version applicative, et des mécanismes
d'orchestration découpant les analyses lourdes en tâches asynchrones. L'efficacité
de la plateforme dépend également de sa capacité à réduire le bruit (faux positifs)
et à transformer les résultats bruts en décisions de remédiation pilotables —
d'où l'importance d'une couche d'enrichissement par intelligence artificielle.

% ─────────────────────────────────────────────────────
\subsection{Besoins fonctionnels}

Les besoins fonctionnels décrivent les capacités que la plateforme doit offrir pour
couvrir l'ensemble du cycle de gestion des vulnérabilités. Chaque fonctionnalité est
motivée par un objectif métier précis, un mécanisme technique sous-jacent et un apport
explicite en termes de sécurité ou de gouvernance.

\begin{itemize}
  \item \textbf{Authentification JWT} : authentification stateless reposant sur des
  jetons d'accès à courte durée de vie et des mécanismes de rafraîchissement, afin
  d'équilibrer sécurité et performance. La contrainte principale porte sur la gestion
  des sessions révoquées et la synchronisation temporelle des services.

  \item \textbf{Authentification forte (2FA TOTP)} : second facteur d'authentification
  basé sur des codes temporels pour les comptes à privilèges, réduisant l'exposition
  aux compromissions par vol d'identifiants.

  \item \textbf{Contrôle d'accès par rôles (RBAC)} : séparation des responsabilités
  via une matrice de droits évaluée à chaque requête, limitant l'accès excessif et
  facilitant l'audit des permissions.

  \item \textbf{Journalisation des actions sensibles (Audit logs)} : enregistrement
  horodaté des opérations critiques — connexion, lancement de scan, changement d'état
  de vulnérabilité — afin de satisfaire aux exigences de conformité et d'investigation.

  \item \textbf{Gestion des projets et des dépôts} : centralisation des projets, dépôts
  Git et URLs cibles, avec paramètres de scan configurables et stockage sécurisé des
  secrets (tokens, clés).

  \item \textbf{Intégration GitHub OAuth} : rattachement des dépôts et branches aux
  projets pour automatiser la collecte de métadonnées et simplifier l'authentification
  des équipes de développement.

  \item \textbf{Clés API} : authentification technique des appels provenant de pipelines
  CI/CD ou d'un CLI, sans exposition d'identifiants utilisateur.

  \item \textbf{Orchestration asynchrone des scans} : planification, exécution en file
  d'attente, suivi de progression et arrêt des analyses, avec gestion de la concurrence
  pour garantir la stabilité opérationnelle en charge élevée.

  \item \textbf{Analyse statique SAST} : détection de patterns de code vulnérables via
  Semgrep et Bandit, avec normalisation des sorties hétérogènes vers un modèle commun.
  La maîtrise des faux positifs constitue l'enjeu principal.

  \item \textbf{Analyse des dépendances (SCA)} : identification des composants tiers
  vulnérables par interrogation de la base OSV, alignée sur les identifiants CVE.

  \item \textbf{Analyse dynamique DAST} : pilotage du moteur OWASP ZAP par API REST,
  avec gestion du contexte applicatif et de l'authentification. La durée d'exécution
  et la disponibilité des environnements cibles constituent les contraintes majeures.

  \item \textbf{Normalisation des vulnérabilités} : harmonisation des sévérités,
  classifications (CWE/OWASP), preuves et localisations issus de sources hétérogènes,
  afin de produire une vision cohérente du risque.

  \item \textbf{Scoring de risque} : calcul d'un score agrégé par projet et par version,
  permettant de prioriser les efforts de remédiation selon des critères objectifs.

  \item \textbf{Workflow de gestion des vulnérabilités} : gestion des états (ouvert,
  confirmé, corrigé, faux positif), assignation à des responsables et suivi par
  commentaires, pour structurer et tracer la remédiation.

  \item \textbf{Enrichissement par IA locale (LLM)} : contextualisation des alertes,
  proposition de correctifs et assistance à la priorisation via l'action
  \textit{Assist with AI}, exécutée sur un modèle local garantissant la confidentialité
  du code.

  \item \textbf{Génération de rapports et exports} : production de livrables PDF et JSON
  filtrables par période, projet et niveau de sévérité, à destination des équipes de
  direction et d'audit.

  \item \textbf{Tableau de bord et visualisation} : consolidation des métriques de
  sécurité, des tendances et des taux de remédiation, avec agrégations optimisées pour
  maintenir la fluidité de l'interface.

  \item \textbf{API REST et intégration CI/CD} : déclenchement et suivi des scans depuis
  les pipelines de livraison, avec contrats d'API stables, gestion des quotas et
  authentification robuste.

  \item \textbf{Notifications et alerting} : diffusion par email, webhooks ou canal
  temps réel lors de la détection de vulnérabilités critiques, avec mécanismes
  anti-saturation configurables.

  \item \textbf{Historisation et traçabilité des scans} : conservation de l'historique
  par projet et par version pour permettre les comparaisons longitudinales, les audits
  et le suivi de la maturité sécurité dans le temps.
\end{itemize}

% ─────────────────────────────────────────────────────
\subsection{Besoins non fonctionnels}

Un besoin non fonctionnel est une condition qui spécifie les critères permettant
de juger le fonctionnement d'un système, plutôt que ses comportements spécifiques.
À la différence des besoins fonctionnels qui définissent \textit{ce que} le système
fait, les besoins non fonctionnels définissent \textit{comment} il le fait, en termes
de performance, de fiabilité, de sécurité et de maintenabilité.

\begin{itemize}
  \item \textbf{Performance et latence} : les opérations de consultation doivent rester
  fluides ; les traitements intensifs (scans, normalisation, IA) sont systématiquement
  déportés vers des workers asynchrones afin de ne pas impacter la réactivité de l'API.

  \item \textbf{Scalabilité} : la montée en charge est assurée par une multiplication
  horizontale des workers d'orchestration, sans dégradation de la qualité de service
  pour les utilisateurs actifs.

  \item \textbf{Disponibilité et continuité de service} : supervision des conteneurs
  et politique de redémarrage automatique pour maintenir la disponibilité dans les
  pipelines CI/CD, même en cas de panne d'un service secondaire.

  \item \textbf{Sécurité et confidentialité} : chiffrement des échanges (TLS),
  protection des secrets (tokens, clés API), cloisonnement des données par projet
  et journalisation exhaustive des accès aux ressources sensibles.

  \item \textbf{Maintenabilité} : architecture modulaire, code documenté et couverture
  de tests suffisante pour faciliter l'évolution des moteurs de scan et des règles
  métier sans régressions.

  \item \textbf{Extensibilité} : intégration de nouveaux scanners ou formats d'alertes
  via une architecture par connecteurs et un modèle de données ouvert, sans remaniement
  du cœur applicatif.

  \item \textbf{Ergonomie et adoption} : parcours utilisateurs cohérents, statuts
  explicites et interface claire, afin de minimiser la friction opérationnelle et
  favoriser l'adoption par les équipes.

  \item \textbf{Traçabilité et auditabilité} : journaux complets incluant l'identité
  de l'acteur, le contexte de l'action et l'horodatage, pour répondre aux exigences
  de conformité et d'investigations post-incident.

  \item \textbf{Portabilité} : déploiement entièrement conteneurisé (Docker Compose),
  garantissant la reproductibilité des environnements et la compatibilité avec les
  principaux gestionnaires de dépôts et systèmes CI/CD du marché.

  \item \textbf{Résilience et fiabilité} : traitement idempotent des tâches et
  mécanismes de reprise automatique, assurant la continuité fonctionnelle en cas
  de pannes partielles (indisponibilité d'un scanner ou du module IA).
\end{itemize}

% ─────────────────────────────────────────────────────
\subsection{Diagramme global des cas d'utilisation}

Les cas d'utilisation représentent les séquences d'actions réalisées par le système
qui produisent un résultat observable pour un acteur donné. Ils permettent de modéliser
les attentes des utilisateurs, de délimiter le périmètre fonctionnel du système et
d'améliorer la compréhension de ses interactions avec l'environnement extérieur.

Le diagramme~\ref{fig:use-case-global} présente les interactions entre les cinq acteurs
identifiés et les principales fonctionnalités d'InvisiThreat. Les relations
\textit{include} matérialisent des étapes systématiquement nécessaires — par exemple,
l'authentification est incluse dans tout déclenchement de scan — tandis que les
relations \textit{extend} représentent des comportements conditionnels, comme
l'export PDF qui étend la consultation des résultats selon le rôle et le contexte.

\begin{figure}[H]
\centering
\fbox{\rule{0pt}{2in}\rule{0.9\linewidth}{0pt}}
\caption{Diagramme de cas d'utilisation global d'InvisiThreat.}
\label{fig:use-case-global}
\end{figure}

% ───────────────────────────────────────────────────────────────────
\section{Gestion du projet avec Scrum}

Le Sprint 0 structure le Product Backlog et prépare la planification des sprints
à venir. La priorisation s'appuie sur la criticité sécurité, les dépendances
techniques et la valeur métier de chaque fonctionnalité.

% ─────────────────────────────────────────────────────
\subsection{Product Backlog}

Le Product Backlog formalise la traduction des besoins en unités de valeur livrables.
Dans une logique Scrum, il organise le périmètre autour d'epics, de features et de
user stories, tout en respectant une priorisation fondée sur le risque et la valeur
attendue.

\subsubsection{Structuration en epics}

Les epics regroupent le backlog par grands domaines fonctionnels. L'epic
\textit{Authentification et gouvernance} couvre la sécurité des accès et la gestion
des rôles. L'epic \textit{Gestion des projets} englobe la configuration des cibles
et des dépôts. L'epic \textit{Scans SAST/DAST/SCA} porte sur l'orchestration
technique des analyses, tandis que l'epic \textit{Résultats, workflow et IA} regroupe
la normalisation, le scoring, la remédiation assistée et les exports. Cette
structuration favorise la lisibilité du backlog et facilite la planification
incrémentale.

\subsubsection{Technique de priorisation du Product Backlog}

La priorisation du Product Backlog conditionne directement la valeur livrée à chaque
sprint. Trois critères complémentaires ont été retenus :

\begin{itemize}
  \item \textbf{Business Value} : priorité aux fonctionnalités apportant le maximum
  de valeur à l'utilisateur final (détection précoce des vulnérabilités, visibilité
  du risque par projet).

  \item \textbf{Risk Minimizing} : traitement anticipé des éléments les plus complexes
  ou les plus incertains sur le plan technique (orchestration asynchrone, intégration
  du module LLM local, pilotage de ZAP par API).

  \item \textbf{MoSCoW} : classification des items selon quatre niveaux —
  \textbf{M} (\textit{Must have}, obligatoire pour la livraison),
  \textbf{S} (\textit{Should have}, fortement souhaitable),
  \textbf{C} (\textit{Could have}, optionnel à valeur ajoutée),
  \textbf{W} (\textit{Would have}, reporté aux versions futures).
\end{itemize}

\subsubsection{Synthèse du Product Backlog}

Les tableaux suivants présentent l'ensemble des user stories issues du cahier des
charges, organisées par domaine fonctionnel et ordonnées selon leur niveau de priorité
MoSCoW.

\begin{table}[H]
\centering
\caption{Product Backlog -- Gouvernance et gestion des accès}
\label{tab:backlog-gouvernance}
\small
\renewcommand{\arraystretch}{1.2}
\begin{tabularx}{\textwidth}{c X c X c}
\hline
\rowcolor{gray!15}
\textbf{ID} & \textbf{Feature} & \textbf{US} & \textbf{User Story} & \textbf{Priorité} \\
\hline
\rowcolor{gray!10}
\multicolumn{5}{c}{\textbf{Module Gouvernance et Accès}} \\
\hline
1 & Authentification & 1.1
  & En tant qu'administrateur, je veux m'authentifier afin d'accéder à la plateforme
    et à l'ensemble des fonctionnalités de gouvernance. & M \\
1 & Authentification & 1.2
  & En tant que développeur, je veux m'authentifier afin d'accéder uniquement à
    mes projets rattachés. & M \\
2 & 2FA TOTP & 2.1
  & En tant qu'administrateur, je veux activer l'authentification à deux facteurs
    afin de renforcer la protection de mon compte. & M \\
3 & RBAC & 3.1
  & En tant qu'administrateur, je veux définir des rôles et des matrices de
    permissions afin de cloisonner les accès selon les responsabilités. & M \\
4 & Audit logs & 4.1
  & En tant que Security Manager, je veux consulter le journal des actions sensibles
    afin d'assurer la traçabilité et la conformité. & M \\
\hline
\rowcolor{gray!10}
\multicolumn{5}{c}{\textbf{Module Projets et Intégrations}} \\
\hline
5 & Gestion des projets & 5.1
  & En tant qu'administrateur, je veux créer un projet et définir ses URLs cibles
    afin de centraliser la configuration des analyses. & M \\
5 & Gestion des projets & 5.2
  & En tant qu'administrateur, je veux configurer les paramètres de scan par projet
    afin d'adapter les analyses au contexte applicatif. & S \\
6 & Gestion des membres & 6.1
  & En tant qu'administrateur, je veux inviter des membres et leur attribuer des rôles
    afin de structurer la collaboration sur la plateforme. & S \\
7 & GitHub OAuth & 7.1
  & En tant que développeur, je veux connecter un dépôt GitHub via OAuth afin
    d'automatiser la collecte des métadonnées de branches. & S \\
8 & Clés API & 8.1
  & En tant que système CI/CD, je veux m'authentifier via une clé API afin de
    déclencher des scans sans exposer d'identifiants utilisateur. & S \\
\hline
\end{tabularx}
\end{table}

\begin{table}[H]
\centering
\caption{Product Backlog -- Scans et orchestration}
\label{tab:backlog-scans}
\small
\renewcommand{\arraystretch}{1.2}
\begin{tabularx}{\textwidth}{c X c X c}
\hline
\rowcolor{gray!15}
\textbf{ID} & \textbf{Feature} & \textbf{US} & \textbf{User Story} & \textbf{Priorité} \\
\hline
\rowcolor{gray!10}
\multicolumn{5}{c}{\textbf{Module Scans SAST / SCA / DAST et Orchestration}} \\
\hline
9 & SAST Semgrep/Bandit & 9.1
  & En tant que développeur, je veux lancer une analyse statique afin de détecter
    les patterns de code vulnérables avant tout déploiement. & M \\
10 & Règles internes & 10.1
  & En tant que Security Manager, je veux appliquer des règles regex personnalisées
    afin d'adapter la détection aux spécificités du code interne. & S \\
11 & SCA OSV & 11.1
  & En tant que développeur, je veux analyser les dépendances via OSV afin
    d'identifier les composants tiers présentant des vulnérabilités connues. & M \\
12 & CLI scanner & 12.1
  & En tant que système CI/CD, je veux déclencher une analyse via la ligne de commande
    afin de l'intégrer directement dans la chaîne de livraison. & S \\
13 & Orchestration asynchrone & 13.1
  & En tant que système, je veux que les scans soient mis en file d'attente via
    Celery/Redis afin de garantir la stabilité sous charge élevée. & M \\
14 & DAST OWASP ZAP & 14.1
  & En tant que développeur, je veux lancer une analyse dynamique via OWASP ZAP
    afin de détecter les failles d'exposition réelles de l'application. & M \\
15 & Normalisation & 15.1
  & En tant que Security Manager, je veux consulter des alertes normalisées selon
    une sévérité unifiée afin de prioriser la remédiation efficacement. & M \\
\hline
\end{tabularx}
\end{table}

\begin{table}[H]
\centering
\caption{Product Backlog -- Workflow, IA, reporting et CI/CD}
\label{tab:backlog-ia-notifications}
\small
\renewcommand{\arraystretch}{1.2}
\begin{tabularx}{\textwidth}{c X c X c}
\hline
\rowcolor{gray!15}
\textbf{ID} & \textbf{Feature} & \textbf{US} & \textbf{User Story} & \textbf{Priorité} \\
\hline
\rowcolor{gray!10}
\multicolumn{5}{c}{\textbf{Module Workflow, IA et Reporting}} \\
\hline
16 & Assistant IA local & 16.1
  & En tant que développeur, je veux obtenir une explication et une recommandation
    de correction via \textit{Assist with AI} afin d'accélérer la remédiation. & M \\
17 & Scoring de risque & 17.1
  & En tant que Security Manager, je veux consulter un score de risque par projet
    afin de prioriser les efforts de sécurité selon des critères objectifs. & M \\
18 & Workflow vulnérabilités & 18.1
  & En tant que développeur, je veux changer l'état d'une vulnérabilité et l'assigner
    à un responsable afin de structurer le processus de remédiation. & M \\
19 & Dashboard global & 19.1
  & En tant que Security Manager, je veux visualiser un tableau de bord consolidé
    afin d'avoir une vue d'ensemble du niveau de risque sur l'ensemble des projets. & S \\
20 & Détail des vulnérabilités & 20.1
  & En tant que développeur, je veux consulter les preuves et le contexte d'une
    alerte afin de comprendre son impact et de la reproduire. & M \\
21 & Génération de rapports & 21.1
  & En tant que Security Manager, je veux générer un rapport PDF ou JSON par projet
    afin de le partager lors des revues de sécurité. & S \\
22 & Historique comparatif & 22.1
  & En tant que Security Manager, je veux comparer les résultats de scan entre
    versions afin de mesurer la progression de la maturité sécurité. & S \\
\hline
\rowcolor{gray!10}
\multicolumn{5}{c}{\textbf{Module CI/CD et Notifications}} \\
\hline
23 & Déclenchement CI/CD & 23.1
  & En tant que système CI/CD, je veux déclencher automatiquement un scan sur
    commit ou merge request afin d'intégrer la sécurité au pipeline de livraison. & M \\
24 & Notifications temps réel & 24.1
  & En tant que développeur, je veux recevoir des notifications instantanées lors
    de la détection de vulnérabilités critiques afin de réagir sans délai. & S \\
25 & Politiques d'alerte & 25.1
  & En tant que Security Manager, je veux configurer des seuils d'alerte par sévérité
    afin d'adapter le niveau de notification au contexte du projet. & S \\
26 & Intégration Slack & 26.1
  & En tant que Security Manager, je veux recevoir les alertes critiques dans un canal
    Slack afin de centraliser les notifications dans l'outil de communication de
    l'équipe. & W \\
27 & Comparaison inter-projets & 27.1
  & En tant que Security Manager, je veux comparer les scores de risque entre projets
    afin d'identifier les périmètres les plus exposés à l'échelle du portefeuille. & W \\
\hline
\end{tabularx}
\end{table}

% ─────────────────────────────────────────────────────
\subsection{Planification des sprints}

La planification de sprint répond à deux questions fondamentales : qu'est-ce qui peut
être livré à l'issue de ce sprint ? Et de quelle façon le travail sélectionné
sera-t-il conduit ? Pour InvisiThreat, la durée totale du projet a été estimée à
\textbf{quatre mois}, découpée en cinq sprints de deux à trois semaines chacun,
conformément aux préconisations Scrum pour des projets de taille intermédiaire.

La répartition des responsabilités au sein de l'équipe est la suivante :

\begin{itemize}
  \item \textbf{Stagiaire 1} : développement backend, module IA/LLM, DevSecOps et
  orchestration des scans.
  \item \textbf{Stagiaire 2} : développement frontend, tableaux de bord et intégration
  CI/CD.
  \item \textbf{Encadrant professionnel} : supervision technique et validation des
  choix d'architecture.
  \item \textbf{Encadrant académique} : suivi méthodologique et validation scientifique.
\end{itemize}

Le planning détaillé est présenté au tableau~\ref{tab:sprints}. La progression suit
une logique de construction en couches : fondations techniques, puis gouvernance des
accès, puis gestion des projets, puis analyses de sécurité, et enfin restitution
enrichie et automatisation.

\begin{table}[H]
\centering
\caption{Planification et livrables des sprints InvisiThreat}
\label{tab:sprints}
\small
\renewcommand{\arraystretch}{1.4}
\begin{tabularx}{\textwidth}{c p{2.8cm} X X p{1.8cm}}
\hline
\rowcolor{gray!15}
\textbf{Sprint} & \textbf{Objectif} & \textbf{Contenu} & \textbf{Livrables}
  & \textbf{Durée} \\
\hline
0 & \textbf{Cadrage et fondations}
  & Initialisation du dépôt, configuration Docker~Compose, modélisation de la
    base de données, mise en place des squelettes FastAPI et React.
  & Socle technique opérationnel, schéma PostgreSQL validé, services démarrables.
  & 2~semaines \\
\hline
1 & \textbf{Sécurité et gouvernance}
  & Authentification JWT, RBAC, 2FA~TOTP, audit logs, politiques de sécurité.
  & Module d'authentification sécurisé et journalisation opérationnelle.
  & 3~semaines \\
\hline
2 & \textbf{Projets et intégrations}
  & CRUD projets et membres, intégration GitHub~OAuth, gestion des clés API,
    notifications de base.
  & Gestion de projets fonctionnelle, intégration GitHub et clés API opérationnelles.
  & 3~semaines \\
\hline
3 & \textbf{SAST et SCA}
  & Intégration Semgrep et Bandit, règles regex internes, analyse OSV, CLI~scanner,
    normalisation des alertes.
  & Pipeline SAST/SCA opérationnel, CLI fonctionnel.
  & 3~semaines \\
\hline
4 & \textbf{DAST, orchestration et IA}
  & Intégration OWASP~ZAP, workers Celery/Redis, suivi temps réel, scoring de risque,
    workflow de vulnérabilités, assistant LLM local.
  & Pipeline DAST asynchrone, scoring de risque, workflow et recommandations IA.
  & 3~semaines \\
\hline
\end{tabularx}
\end{table}

% ───────────────────────────────────────────────────────────────────
\section{Étude technique}

\subsection{Environnement matériel}

Le projet a été réalisé sur deux postes de travail homogènes, un par stagiaire, afin
de garantir des conditions de développement et de test reproductibles. Ce dimensionnement
répond à la nature des traitements effectués : l'analyse statique nécessite un accès
rapide au système de fichiers et une charge CPU stable, tandis que l'analyse dynamique
sollicite davantage la mémoire et la latence réseau pour piloter les scans actifs.
L'isolation des conteneurs — notamment OWASP ZAP — dans des processus dédiés limite
la contention avec l'API et l'interface, et contribue à stabiliser les temps de réponse
observés lors des tests.

\begin{table}[H]
\centering
\caption{Configuration matérielle des postes de travail}
\label{tab:env-materiel}
\renewcommand{\arraystretch}{1.2}
\begin{tabular}{|l|c|c|}
\hline
\rowcolor{gray!15}
\textbf{Élément} & \textbf{PC 1 -- Wiem Bziouich} & \textbf{PC 2 -- Nouha Boughnim} \\
\hline
Processeur     & Intel Core i7 & Intel Core i7 \\
RAM installée  & 16 Go         & 16 Go         \\
Système        & Windows 11 Pro 64 bits & Windows 11 Pro 64 bits \\
Stockage       & SSD 512 Go    & SSD 512 Go    \\
\hline
\end{tabular}
\end{table}

% ─────────────────────────────────────────────────────
\subsection{Environnement logiciel}

Le tableau~\ref{tab:env-logiciel} regroupe l'ensemble des outils utilisés par l'équipe
pour le développement, les tests, la modélisation et la gestion du projet. Ces outils
constituent l'environnement de travail au sens ingénierie logicielle, distinct de la
stack technique déployée en production.

% Logos helpers for consistent size and alignment
\newcommand{\logocell}[1]{\raisebox{-0.5\height}{\includegraphics[width=0.75in,height=0.45in,keepaspectratio]{#1}}}
\newcommand{\placeholderlogo}{\fbox{\rule{0pt}{0.45in}\rule{0.75in}{0pt}}}

\begin{table}[H]
\centering
\caption{Environnement logiciel utilisé pour le développement}
\label{tab:env-logiciel}
\renewcommand{\arraystretch}{2.0}
\begin{tabularx}{\textwidth}{c p{3cm} X}
\hline
\rowcolor{gray!15}
\textbf{Logo} & \textbf{Outil} & \textbf{Description et rôle dans le projet} \\
\hline
\logocell{images/logos/VSCode.png}
  & \textbf{VS Code}
  & IDE principal pour le développement Python et JavaScript, avec extensions de
    linting, d'intégration Git et de test des APIs REST. \\
\hline
\logocell{images/logos/Git.png}
  & \textbf{Git / GitHub}
  & Contrôle de version distribué, collaboration par branches et revue de code pour
    assurer la traçabilité des évolutions. \\
\hline
\logocell{images/logos/Docker.png}
  & \textbf{Docker Desktop}
  & Exécution locale des conteneurs pour reproduire fidèlement l'environnement
    de production et éviter les conflits de dépendances. \\
\hline
\placeholderlogo
  & \textbf{Postman}
  & Test des endpoints REST, validation des payloads et reproduction des scénarios
    de scan pour vérifier les contrats d'API. \\
\hline
\logocell{images/logos/Draw.png}
  & \textbf{Draw.io}
  & Modélisation des diagrammes UML : cas d'utilisation, déploiement et architecture
    logique. \\
\hline
\placeholderlogo
  & \textbf{Trello / Jira / ClickUp}
  & Gestion du Product Backlog, suivi de l'avancement des sprints et tableaux Kanban
    partagés entre les membres de l'équipe. \\
\hline
\placeholderlogo
  & \textbf{Chrome / Firefox DevTools}
  & Validation de l'interface React, inspection des requêtes réseau et analyse des
    performances côté client. \\
\hline
\placeholderlogo
  & \textbf{venv / pip}
  & Isolation des dépendances Python par environnement virtuel, garantissant la
    cohérence entre les deux postes de développement. \\
\hline
\logocell{images/logos/node.png}
  & \textbf{Node.js / npm}
  & Gestion des dépendances et build du frontend React, avec contrôle précis des
    versions des bibliothèques côté client. \\
\hline
\end{tabularx}
\end{table}

% ─────────────────────────────────────────────────────
\subsection{Technologies et outils utilisés}

Le tableau~\ref{tab:technologies} présente la stack technique déployée pour assurer
la détection, la normalisation et la restitution des résultats de sécurité.
Chaque choix est justifié par des exigences de performance, de sécurité et de
maintenabilité, et accompagné d'une mise en perspective par rapport aux alternatives.

\begin{table}[H]
\centering
\caption{Stack technique d'InvisiThreat -- technologies et justifications}
\label{tab:technologies}
\renewcommand{\arraystretch}{2.0}
\begin{tabularx}{\textwidth}{c p{3cm} X}
\hline
\rowcolor{gray!15}
\textbf{Logo} & \textbf{Technologie} & \textbf{Rôle, justification et contraintes} \\
\hline
\logocell{images/logos/python.png}
  & \textbf{Python / FastAPI} \newline \small{Backend API}
  & Expose une API REST asynchrone capable de traiter des requêtes concurrentes
    sans blocage lors des scans longs. FastAPI apporte un typage strict et une
    validation automatique des schémas (Pydantic), ce qui réduit les erreurs
    d'intégration par rapport à Flask. Sa structure légère est mieux adaptée
    qu'un framework Django pour un service d'orchestration découplé. \\
\hline
\logocell{images/logos/react.png}
  & \textbf{React} \newline \small{Frontend}
  & Bibliothèque de construction d'interfaces riches et réactives, adaptée à la
    visualisation temps réel des états de scan. Son écosystème étendu et sa
    flexibilité de composition des composants le distinguent de Vue ou Angular
    pour des tableaux de bord complexes. La gestion d'état est assurée via
    hooks et contextes React. \\
\hline
\logocell{images/logos/postgreSQL.png}
  & \textbf{PostgreSQL} \newline \small{Base de données}
  & Assure la persistance transactionnelle (ACID) et les agrégations avancées
    nécessaires à l'historisation des scans et aux analyses de tendances par
    sévérité. Offre une richesse fonctionnelle supérieure à MySQL pour les
    requêtes analytiques. La contrainte principale est l'optimisation des index
    sur de grands volumes d'alertes. \\
\hline
\logocell{images/logos/redis.png}
  & \textbf{Redis} \newline \small{Broker \& cache}
  & Utilisé comme broker de messages pour Celery et stockage temporaire des
    états de scan (progression, files d'attente, verrous distribués). Sa nature
    en mémoire garantit des temps d'accès très faibles et décharge PostgreSQL
    pour les données éphémères. La stratégie d'expiration des clés constitue
    la contrainte principale. \\
\hline
\placeholderlogo
  & \textbf{Celery} \newline \small{Orchestration async.}
  & Exécute les analyses SAST, SCA et DAST dans des workers isolés sans bloquer
    l'API principale. Fournit un mécanisme de file d'attente, de retries
    configurables et de suivi d'exécution adapté aux tâches longues. La
    supervision des workers et la gestion fine de la concurrence constituent
    les contraintes opérationnelles majeures. \\
\hline
\logocell{images/logos/Socket.IO.png}
  & \textbf{Socket.IO} \newline \small{Temps réel}
  & Assure la diffusion en temps réel des états de scan et des notifications
    critiques vers le frontend, éliminant la latence d'une interrogation
    périodique (polling). La gestion des connexions persistantes et la montée
    en charge représentent les contraintes principales. \\
\hline
\logocell{images/logos/Docker.png}
  & \textbf{Docker / Docker Compose} \newline \small{Conteneurisation}
  & Constitue le socle de déploiement de l'ensemble des services, garantissant
    la reproductibilité des environnements et la séparation des composants
    critiques. Réduit la surcharge par rapport aux machines virtuelles et
    s'intègre naturellement aux pratiques DevSecOps. La gouvernance des images
    et la sécurisation des volumes sont les contraintes clés. \\
\hline
\logocell{images/logos/owaspzap.png}
  & \textbf{OWASP ZAP} \newline \small{DAST}
  & Moteur d'analyse dynamique open source retenu pour son API REST native
    permettant un pilotage programmatique complet. Plus accessible et
    automatisable dans un pipeline CI/CD que Burp Suite Professional, au prix
    d'une précision légèrement inférieure sur certaines règles, compensée par
    un paramétrage fin des profils de scan. \\
\hline
\placeholderlogo
  & \textbf{Semgrep / Bandit} \newline \small{SAST}
  & Moteurs d'analyse statique complémentaires : Semgrep pour la détection
    polyglotte par règles configurables, Bandit pour l'inspection spécialisée
    du code Python. La maîtrise du taux de faux positifs et la maintenance
    des règles constituent les enjeux principaux. \\
\hline
\placeholderlogo
  & \textbf{OSV} \newline \small{SCA}
  & Base de données ouverte de vulnérabilités maintenue collaborativement,
    permettant d'identifier les composants tiers présentant des failles connues
    alignées sur les identifiants CVE. La fréquence d'actualisation des
    métadonnées et la gestion des faux positifs sont les contraintes
    principales. \\
\hline
\placeholderlogo
  & \textbf{Ollama (Mistral / Qwen)} \newline \small{Assistant IA local}
  & Exécute localement des modèles LLM pour contextualiser les alertes et
    proposer des recommandations de remédiation via l'action
    \textit{Assist with AI}, garantissant la confidentialité totale du code
    source. Les ressources matérielles nécessaires à l'inférence et la qualité
    variable des réponses selon le modèle retenu constituent les contraintes
    principales. \\
\hline
\placeholderlogo
  & \textbf{Pytest} \newline \small{Tests automatisés}
  & Framework de tests unitaires et d'intégration pour le backend, offrant un
    modèle de fixtures flexible et des assertions lisibles, supérieur au module
    standard \texttt{unittest}. La maintenance des suites de tests lors des
    évolutions rapides du schéma d'API constitue la contrainte principale. \\
\hline
\end{tabularx}
\end{table}

% ─────────────────────────────────────────────────────
\subsection{Architecture du système}

\subsubsection{Architecture logicielle}

L'architecture logicielle d'InvisiThreat repose sur une séparation claire des
responsabilités en quatre couches fonctionnelles :

\begin{itemize}
  \item \textbf{Couche API} : gestion des projets, des accès utilisateurs, du
  déclenchement des scans et de l'exposition des résultats via des endpoints REST
  documentés.
  \item \textbf{Couche d'orchestration} : exécution asynchrone des analyses longues
  (SAST, SCA, DAST) via des workers Celery, avec file d'attente Redis et suivi
  de progression.
  \item \textbf{Couche de persistance} : stockage normalisé des vulnérabilités,
  de l'historique des scans et des artefacts d'audit dans PostgreSQL.
  \item \textbf{Couche de présentation} : interface React pour la visualisation,
  le workflow de remédiation et les notifications temps réel via Socket.IO.
\end{itemize}

Le flux principal suit une chaîne continue : déclenchement d'un scan, exécution
asynchrone via les workers, normalisation des alertes vers le modèle commun,
calcul du scoring de risque, enrichissement par le module IA, puis exposition
des résultats via les endpoints REST et le canal temps réel.

\begin{figure}[H]
\centering
\fbox{\rule{0pt}{2in}\rule{0.9\linewidth}{0pt}}
\caption{Architecture logicielle en couches d'InvisiThreat.}
\label{fig:archi-logicielle}
\end{figure}

\subsubsection{Architecture physique}

Sur le plan physique, InvisiThreat est déployé via \textbf{Docker Compose} selon
une organisation en trois niveaux :

\begin{itemize}
  \item \textbf{Niveau client} : navigateur React servi par un conteneur Nginx,
  communiquant avec le backend via HTTP/REST et Socket.IO.
  \item \textbf{Niveau serveur d'application} : conteneurs FastAPI, workers Celery,
  broker Redis et module IA Ollama, isolés et communicants par réseau Docker interne.
  \item \textbf{Niveau données} : conteneur PostgreSQL pour la persistance, avec
  volumes montés pour la durabilité des données entre redémarrages.
\end{itemize}

InvisiThreat adopte une architecture de \textbf{monolithe modulaire} côté backend,
enrichie par une orchestration asynchrone. Les moteurs de scan (ZAP, Semgrep, Bandit)
et le module IA sont isolés dans des conteneurs dédiés, sans imposer la complexité
opérationnelle d'une architecture microservices complète. La scalabilité horizontale
est assurée par la multiplication des workers Celery ; la tolérance aux pannes est
renforcée car la défaillance d'un scanner n'interrompt pas les services de consultation ;
et l'observabilité (logs structurés, métriques, audit trail) garantit une traçabilité
de bout en bout.

\begin{figure}[H]
\centering
\fbox{\rule{0pt}{2in}\rule{0.9\linewidth}{0pt}}
\caption{Architecture physique -- déploiement Docker Compose d'InvisiThreat.}
\label{fig:archi-physique}
\end{figure}

% ───────────────────────────────────────────────────────────────────
\section{Diagramme de déploiement}

Le diagramme de déploiement illustre l'organisation conteneurisée de la plateforme
via Docker Compose, les flux réseau entre services et l'exposition des ports
applicatifs (API FastAPI, frontend React, moteur ZAP, module Ollama). Il permet
de visualiser l'isolation des composants critiques et les dépendances de
communication entre niveaux.

\begin{figure}[H]
\centering
\fbox{\rule{0pt}{2in}\rule{0.9\linewidth}{0pt}}
\caption{Diagramme de déploiement Docker Compose d'InvisiThreat.}
\label{fig:comp-deployment}
\end{figure}

% ───────────────────────────────────────────────────────────────────
\section{Conclusion}

Ce chapitre a posé les fondations méthodologiques et techniques du projet InvisiThreat.
L'identification des acteurs et la spécification exhaustive des besoins fonctionnels
et non fonctionnels ont permis de délimiter précisément le périmètre de la plateforme.
La structuration du Product Backlog selon une priorisation MoSCoW et la planification
en cinq sprints garantissent une progression maîtrisée, alignée sur la valeur métier
et la gestion du risque technique. Enfin, l'étude technique a confirmé la cohérence
de la stack retenue avec les exigences DevSecOps : architecture modulaire,
orchestration asynchrone, enrichissement par IA locale et déploiement intégralement
conteneurisé.

Le chapitre suivant présente le \textbf{Sprint 1}, consacré à la sécurité et à
la gouvernance : authentification JWT, contrôle d'accès par rôles, double facteur
TOTP et journalisation des actions sensibles.